\chapter{Evaluation and Discussion}
\label{chap:eval}

\section{Overview}
\label{sec:eval_overview}

This chapter presents the evaluation of the architecture-aware LLM-based maintainability improvement workflow investigated in this thesis. The evaluation was conducted through three sequential empirical studies. Each subsequent study was motivated by limitations or open questions identified in the preceding one.

\textbf{Study 1} executed the full architecture-aware pipeline on 162 Python backend repositories. The pipeline detected repository architecture, selected an architectural tactic, implemented code modifications, and re-measured maintainability. This study revealed that LLM-implemented tactics can produce statistically significant but modest improvements, and that architecture detection is a critical but insufficiently validated pipeline stage.

\textbf{Study 2} was motivated by the lack of independent validation of the architecture detection stage in Study 1. The architecture detection component was evaluated separately using a manually labeled dataset of 57 repositories, five LLMs, and four evidence configurations. This study provided clear evidence about which inputs are useful, which are harmful, and where systematic errors occur.

\textbf{Study 3} applied architectural tactic implementation to the architecture-labeled 57-repository dataset. With ground-truth architecture labels available, this study tested whether more reliable architectural context produces better tactic implementation outcomes. The results were limited.

The three studies use different datasets because they answer different questions. Study 1 requires many repositories to assess pipeline robustness and aggregate maintainability impact. Study 2 requires manually labeled ground truth to evaluate classification accuracy. Study 3 uses the ground-truth labeled dataset to isolate the effect of architecture label quality on implementation.

\subsection{Research Questions}
\label{sec:eval_rq}

The evaluation is organized around three research questions derived from the thesis research objectives:

\begin{itemize}
    \item \textbf{RQ1:} Which LLMs and repository evidence types enable reliable architectural style classification of Python backend repositories?
    \item \textbf{RQ2:} To what extent does an LLM-implemented architectural tactic pipeline produce measurable maintainability improvements?
    \item \textbf{RQ3:} Which architectural styles and tactics are most amenable to LLM-driven maintainability improvement?
\end{itemize}

RQ1 is addressed in Study 2. RQ2 is primarily addressed in Study 1, and partially revisited in Study 3. RQ3 is addressed through both Study 1 and Study 3.

\section{Study 1: Pipeline-Level Maintainability Evaluation}
\label{sec:eval_pipeline}

The first study evaluated the full architecture-aware maintainability improvement pipeline on 162 open-source Python backend repositories.

\subsection{Pipeline Completion and Failure Rate}
\label{subsec:eval_pipeline_completion}

Out of 162 repositories, the pipeline produced paired before/after maintainability measurements for 121 repositories. The remaining 41 repositories produced null outcomes because the LLM introduced modifications that resulted in syntax errors, broken imports, or other issues preventing re-analysis. The overall failure rate was 25.3\%.

Table~\ref{tab:eval_outcomes} shows the distribution of pipeline outcomes.

\begin{table}[H]
\centering
\small
\caption{Pipeline outcome distribution across 162 repositories.}
\label{tab:eval_outcomes}
\begin{tabular}{lrr}
\toprule
\textbf{Outcome} & \textbf{N} & \textbf{Proportion} \\
\midrule
Improved ($\Delta MI > 0$) & 22 & 13.6\% \\
Stable ($\Delta MI = 0$) & 92 & 56.8\% \\
Degraded ($\Delta MI < 0$) & 7 & 4.3\% \\
Failed / null & 41 & 25.3\% \\
\midrule
\textbf{Total} & \textbf{162} & \textbf{100\%} \\
\bottomrule
\end{tabular}
\end{table}

The dominant outcome was stability: 92 of 121 completed repositories showed no measurable Maintainability Index change. This does not mean that no code was modified. Inspection of preserved modification logs showed that even MI-stable repositories had an average of 4.7 code modification steps. However, many modifications affected file-level complexity in ways that did not propagate to MI, or introduced compensating changes that offset each other.

\subsection{Maintainability Index Impact}
\label{subsec:eval_mi_impact}

Table~\ref{tab:eval_mi_descriptive} summarizes the Maintainability Index before and after intervention across the 121 completed repositories.

\begin{table}[H]
\centering
\small
\caption{Maintainability Index descriptive statistics before and after LLM-implemented tactic application (N = 121).}
\label{tab:eval_mi_descriptive}
\begin{tabular}{lrr}
\toprule
\textbf{Metric} & \textbf{MI before} & \textbf{MI after} \\
\midrule
Mean & 65.11 & 65.59 \\
Median & 68.42 & 68.67 \\
Standard deviation & 19.17 & 19.22 \\
\midrule
Mean $\Delta MI$ & \multicolumn{2}{c}{+0.48} \\
Median $\Delta MI$ & \multicolumn{2}{c}{0.00} \\
\bottomrule
\end{tabular}
\end{table}

A Shapiro-Wilk test confirmed non-normality of the delta distribution ($W = 0.34$, $p < 0.001$), justifying the use of a non-parametric test. A Wilcoxon signed-rank test on all 121 paired observations produced:

\begin{equation}
W = 360.0,\quad p = 0.001,\quad r = 0.28
\end{equation}

The result is statistically significant, but the effect size is small ($r = 0.28$). The median $\Delta MI$ of 0.00 indicates that improvements are concentrated in a subset of repositories rather than distributed uniformly.

Among the 22 repositories that improved, the mean improvement was:
\begin{equation}
\overline{\Delta MI}_{improved} = +2.89
\end{equation}

Individual improvements included:
\begin{itemize}
    \item \textit{get\_subscribe}: $\Delta MI = +13.52$ (Decomposability tactic);
    \item \textit{ps4-exploit-host}: $\Delta MI = +11.45$ (Decomposability tactic);
    \item \textit{DXY-COVID-19-Crawler}: $\Delta MI = +9.99$ (Decomposability tactic).
\end{itemize}

These results address \textbf{RQ2}: the pipeline produces a statistically significant but modest overall maintainability improvement, concentrated in a subset of repositories where structure is low and decomposition is feasible.

\subsection{Results by Architectural Tactic}
\label{subsec:eval_tactic_results}

Table~\ref{tab:eval_tactics} shows pipeline outcomes grouped by the architectural tactic selected by the LLM.

\begin{table}[H]
\centering
\small
\caption{Pipeline results grouped by selected architectural tactic.}
\label{tab:eval_tactics}
\begin{tabular}{lrrrrrr}
\toprule
\textbf{Tactic} & \textbf{N} & \textbf{Null} & \textbf{Improved} & \textbf{Stable} & \textbf{Degraded} & $\overline{\Delta MI}$ \\
\midrule
Decomposability & 103 & 29 & 7 & 63 & 4 & +0.55 \\
Localized Modification & 32 & 5 & 6 & 19 & 2 & +0.13 \\
Reduced Coupling & 23 & 5 & 8 & 9 & 1 & +0.70 \\
Other & 4 & 2 & 1 & 1 & 0 & +0.67 \\
\bottomrule
\end{tabular}
\end{table}

Decomposability was the most frequently selected tactic, applied in 103 of 162 repositories. This reflects the composition of the dataset: most repositories exhibited monolithic or loosely structured code, which makes decomposition a natural first-choice tactic.

Reduced Coupling showed the highest improvement rate: 8 of 18 completed cases improved, giving an improvement rate of 44.4\%. This suggests that coupling-reduction transformations can be effective when the LLM can identify and modify dependency relationships without breaking execution.

Decomposability also had the highest null rate among completed cases. Decomposing unstructured code is technically more complex than introducing abstractions, which may explain why Decomposability produces both the largest gains and the most failures.

\subsection{Results by Architectural Style}
\label{subsec:eval_arch_results}

Table~\ref{tab:eval_arch_results} shows pipeline outcomes grouped by the architectural style detected by the LLM.

\begin{table}[H]
\centering
\small
\caption{Pipeline results grouped by detected architectural style.}
\label{tab:eval_arch_results}
\begin{tabular}{lrrrrr}
\toprule
\textbf{Architecture} & \textbf{N} & \textbf{Null \%} & \textbf{Improved} & \textbf{Degraded} & $\overline{\Delta MI}$ \\
\midrule
Script-based & 71 & 29.6 & 5 & 0 & +0.81 \\
Modular monolith & 83 & 19.3 & 15 & 7 & +0.22 \\
Layered & 6 & 50.0 & 2 & 0 & +0.94 \\
MVC & 2 & 50.0 & 0 & 0 & 0.00 \\
\bottomrule
\end{tabular}
\end{table}

Script-based repositories produced the highest mean improvement of +0.81 and contained all three largest individual gains. Script-based repositories often consist of large files with high Halstead Volume, which is directly reduced when files are split into cohesive modules.

Modular monoliths showed more modest mean improvement (+0.22) and accounted for all seven degradation cases. This asymmetry indicates that already-structured systems are sensitive to poorly targeted modifications: the LLM may inadvertently increase coupling or break internal boundaries in code that was already organized.

Layered repositories showed a high mean improvement (+0.94) with no degradations, but the sample was too small (three completed repositories) for reliable conclusions.

These results partially address \textbf{RQ3}: script-based repositories with decomposable monolithic files represent the most receptive targets for the current pipeline.

\subsection{Key Limitation Revealed by Study 1}
\label{subsec:eval_pipeline_limitation}

The initial pipeline experiment relied on LLM-reported confidence (median 0.95, range 0.85--0.98) as the only quality indicator for architecture detection. High concentration of certain tactics in specific architecture types suggested that detection may have been working, but this could not be verified without independent ground truth.

The 25.3\% failure rate and the prevalence of neutral outcomes also suggested that implementation reliability was insufficient and that architecture detection errors may propagate into tactic selection. These limitations motivated the focused architecture detection study in Study 2.

\section{Study 2: Architecture Detection Evaluation}
\label{sec:eval_arch_detection}

Study 2 was conducted after Study 1 revealed that architecture detection is a critical but unvalidated stage of the pipeline. The study evaluated five LLMs and four context configurations on a manually labeled dataset of 57 Python backend repositories.

\subsection{Overall Classification Performance}
\label{subsec:eval_arch_overall}

Table~\ref{tab:eval_arch_accuracy} shows accuracy and macro-F1 for all model-prompt combinations. The majority-class baseline, which always predicts \textit{modular monolith}, yields 57.9\%.

\begin{table}[H]
\centering
\small
\caption{Architecture detection accuracy and macro-F1 across models and prompt configurations. Best result in bold.}
\label{tab:eval_arch_accuracy}
\begin{tabular}{l|cc|cc|cc|cc}
\toprule
 & \multicolumn{2}{c|}{\textbf{P1}} & \multicolumn{2}{c|}{\textbf{P2}} & \multicolumn{2}{c|}{\textbf{P3}} & \multicolumn{2}{c}{\textbf{P4}} \\
\textbf{Model} & Acc & F1$_m$ & Acc & F1$_m$ & Acc & F1$_m$ & Acc & F1$_m$ \\
\midrule
Qwen3    & .596 & .53 & \textbf{.702} & \textbf{.65} & .632 & .52 & .667 & .57 \\
MiniMax  & .579 & .55 & .649 & .61 & .649 & .55 & .649 & .56 \\
Gemma4   & .561 & .48 & .632 & .55 & .614 & .47 & .614 & .45 \\
DeepSeek & .544 & .46 & .614 & .50 & .614 & .48 & .632 & .49 \\
Gemini   & .526 & .47 & .579 & .49 & .561 & .41 & .596 & .45 \\
\bottomrule
\end{tabular}
\end{table}

The best configuration was Qwen3 with P2 (file tree + import graph): 70.2\% accuracy, 0.65 macro-F1. This exceeds the majority-class baseline by 12.3 percentage points and demonstrates that LLM-based architecture detection can provide useful automated triage. However, it is not reliable enough to replace expert judgment, especially at architecture boundaries.

The most important pattern in the results is not model ranking but the consistent effect of evidence type: the delta between P1 and P2 is positive for all five models, while the delta between P2 and P3 is zero or negative for all five models.

This addresses \textbf{RQ1 (partial)}: LLMs can classify repository architecture at above-baseline accuracy when provided with file trees and import dependency graphs.

\subsection{Impact of Repository Evidence Type}
\label{subsec:eval_evidence}

Table~\ref{tab:eval_evidence_delta} shows the accuracy change when each evidence type is added.

\begin{table}[H]
\centering
\small
\caption{Accuracy change when each evidence type is added to the prompt.}
\label{tab:eval_evidence_delta}
\begin{tabular}{lrrr}
\toprule
\textbf{Model} & \textbf{P1 $\to$ P2} & \textbf{P2 $\to$ P3} & \textbf{P3 $\to$ P4} \\
 & \textbf{(+import graph)} & \textbf{(+signatures)} & \textbf{(+metrics)} \\
\midrule
Qwen3    & +10.5\% & $-$7.0\% & +3.5\% \\
MiniMax  & +7.0\%  &  0.0\%  &  0.0\% \\
Gemma4   & +7.0\%  & $-$1.8\% &  0.0\% \\
DeepSeek & +7.0\%  &  0.0\%  & +1.8\% \\
Gemini   & +5.3\%  & $-$1.8\% & +3.5\% \\
\bottomrule
\end{tabular}
\end{table}

Adding the import dependency graph improved accuracy for every model. The import graph encodes inter-module dependencies, which directly reflect how a repository is structurally organized. File trees capture only naming hierarchy; import graphs expose coupling direction and internal modular boundaries.

Adding code signatures degraded or did not change performance for all models. Code signatures introduce a large volume of naming information that may bias the model toward surface-level patterns such as class names containing \texttt{Service}, \texttt{Repository}, or \texttt{Controller}, which can be misleading when distinguishing between modular monolith and layered styles.

Adding structural metrics provided marginal and inconsistent gains.

The conclusion is unambiguous: the optimal evidence configuration for LLM-based architecture detection is file tree combined with import dependency graph. Code signatures should be omitted.

This fully addresses \textbf{RQ1}: import dependency graphs are the most valuable evidence type, while code signatures actively harm performance.

\subsection{Per-Class Classification Performance}
\label{subsec:eval_per_class}

Table~\ref{tab:eval_per_class} shows per-class accuracy at P2.

\begin{table}[H]
\centering
\small
\caption{Per-class architecture detection accuracy at P2 (file tree + import graph).}
\label{tab:eval_per_class}
\begin{tabular}{lccccc}
\toprule
\textbf{Class (N)} & \textbf{Qwen3} & \textbf{MiniMax} & \textbf{Gemma4} & \textbf{DeepSeek} & \textbf{Gemini} \\
\midrule
Modular monolith (33) & \textbf{.727} & .636 & .606 & .515 & .424 \\
Layered (16)          & .750 & .688 & .750 & .750 & .750 \\
Script-based (8)      & .750 & .750 & .750 & \textbf{.875} & .750 \\
\bottomrule
\end{tabular}
\end{table}

The script-based and layered styles are detected reliably across all models, with minimum accuracy of 0.688. The modular monolith class is the primary source of classification errors, with accuracy ranging from 0.424 (Gemini) to 0.727 (Qwen3). Since modular monolith accounts for 57.9\% of the dataset, errors in this class drive most of the overall accuracy differences between models.

\subsection{Systematic Misclassification Patterns}
\label{subsec:eval_misclassification}

Table~\ref{tab:eval_mm_errors} shows how modular monolith repositories were misclassified at P2.

\begin{table}[H]
\centering
\small
\caption{Modular monolith misclassifications at P2 (N = 33 per model).}
\label{tab:eval_mm_errors}
\begin{tabular}{lrrrr}
\toprule
\textbf{Model} & \textbf{Correct} & $\to$ \textbf{layered} & $\to$ \textbf{script} & $\to$ \textbf{other} \\
\midrule
Qwen3    & 24 & 8  & 1 & 0 \\
MiniMax  & 21 & 7  & 2 & 3 \\
Gemma4   & 20 & 9  & 3 & 1 \\
DeepSeek & 17 & 11 & 5 & 0 \\
Gemini   & 14 & 14 & 4 & 1 \\
\bottomrule
\end{tabular}
\end{table}

The dominant error is modular monolith being classified as layered. This is a systematic pattern: all five models exhibit it, and it does not diminish significantly with richer input. The distinction between modular monolith and layered requires understanding whether packages represent technical responsibilities (presentation, service, data) or domain features (users, orders, products). This distinction is semantic and depends on naming intent, which is not fully visible from import topology alone.

\subsection{Confidence Calibration}
\label{subsec:eval_confidence}

Table~\ref{tab:eval_confidence} compares mean confidence for correct and incorrect predictions at P2.

\begin{table}[H]
\centering
\small
\caption{Mean LLM confidence for correct vs. incorrect predictions at P2.}
\label{tab:eval_confidence}
\begin{tabular}{lccc}
\toprule
\textbf{Model} & \textbf{Correct} & \textbf{Incorrect} & $\Delta$ \\
\midrule
MiniMax  & .838 & .812 & .026 \\
Qwen3    & .893 & .876 & .017 \\
Gemma4   & .918 & .904 & .014 \\
DeepSeek & .853 & .839 & .014 \\
Gemini   & .916 & .907 & .009 \\
\bottomrule
\end{tabular}
\end{table}

The maximum observed difference between correct and incorrect confidence is 0.026, which is too small for practical filtering. Gemini reports the highest mean confidence despite having the lowest accuracy among the five models --- a clear calibration inversion. This finding has a direct implication for the pipeline used in Study 1: the self-reported confidence values used to assess architecture detection quality in that experiment cannot be treated as accuracy indicators.

\section{Study 3: Tactic Implementation on Architecture-Labeled Repositories}
\label{sec:eval_study3}

Study 3 was motivated by the question of whether providing the pipeline with better-quality architecture information would improve implementation outcomes. After Study 2 produced manually validated architecture labels for 57 repositories, those labels were used as input to the tactic implementation stage, replacing the LLM's self-reported classifications.

\subsection{Tactic Implementation on Architecture-Labeled Repositories}
\label{subsec:eval_tactic_impl}

This evaluation applies the tactic selection and implementation pipeline to the 57 manually labeled repositories from the architecture detection study, using validated architecture labels as input. The goal is to isolate tactic implementation quality from architecture detection errors.

\subsubsection*{Outcomes}

Of 56 processable repositories, 14 (25.0\%) failed during cloning. The remaining 42 produced paired measurements. Table~\ref{tab:exp5_outcomes} summarizes the outcome distribution and Table~\ref{tab:exp5_all} the aggregate statistics.

\begin{table}[H]
\centering
\small
\caption{Pipeline outcome distribution and MI statistics (N = 56).}
\label{tab:exp5_outcomes}
\begin{tabular}{lrr}
\toprule
\textbf{Outcome} & \textbf{N} & \textbf{Share of completed} \\
\midrule
Improved ($\Delta MI > 0$)  & 18 & 42.9\% \\
Stable ($\Delta MI = 0$)    & 17 & 40.5\% \\
Worsened ($\Delta MI < 0$)  & 7  & 16.7\% \\
Failed / null               & 14 & 25.0\% of all \\
\bottomrule
\end{tabular}
\end{table}

\begin{table}[H]
\centering
\small
\caption{$\Delta MI$ statistics for completed repositories (N = 42).}
\label{tab:exp5_all}
\begin{tabular}{lr}
\toprule
\textbf{Metric} & \textbf{Value} \\
\midrule
Mean $\Delta MI$ (all completed)  & +1.484 \\
Mean $\Delta MI$ (improved only)  & +3.716 \\
Mean $\Delta MI$ (worsened only)  & $-$0.653 \\
\bottomrule
\end{tabular}
\end{table}

Compared to the initial pipeline experiment (Section~\ref{sec:eval_pipeline}), both the improvement rate (42.9\% vs.\ 18.2\%) and the mean gain (+1.484 vs.\ +0.48) are substantially higher. This is consistent with the hypothesis that validated architecture labels produce more coherent tactic selection.

\subsubsection*{Results by Tactic, Architecture, and Repository Size}

Table~\ref{tab:exp5_combined} shows results broken down by the three primary grouping variables.

\begin{table}[H]
\centering
\small
\caption{Results by tactic, architectural style, and repository size.}
\label{tab:exp5_combined}
\begin{tabular}{llrrrl}
\toprule
\textbf{Group} & \textbf{Category} & \textbf{N} & \textbf{Impr.} & \textbf{Worse.} & $\overline{\Delta MI}$ \\
\midrule
\multirow{4}{*}{Tactic}
  & Decomposability        & 18 & 8 & 3 & +2.33 \\
  & Localized Modification & 14 & 9 & 2 & +1.49 \\
  & Reduced Coupling       & 3  & 1 & 2 & $-$0.16 \\
  & (none)                 & 7  & 0 & 0 & 0.00 \\
\midrule
\multirow{3}{*}{Architecture}
  & Script-based     & 7  & 3 & 0 & +5.62 \\
  & Modular monolith & 26 & 9 & 5 & +0.76 \\
  & Layered          & 9  & 6 & 2 & +0.35 \\
\midrule
\multirow{4}{*}{Size}
  & Tiny ($<$10 files)    & 11 & — & — & median $|\Delta|$ = 1.04 \\
  & Small (10--30 files)  & 9  & — & — & median $|\Delta|$ = 0.76 \\
  & Medium (30--80 files) & 12 & — & — & median $|\Delta|$ = 0.05 \\
  & Large ($>$80 files)   & 10 & — & — & median $|\Delta|$ = 0.00 \\
\bottomrule
\end{tabular}
\end{table}

\textbf{Tactic.} Localized Modification achieves the best risk-adjusted performance (9 improvements, 2 worsenings), while Decomposability is high-variance: it produces the largest gains (+21.98) but also generates more worsenings (3 cases) when new extracted modules have MI below the repository average.

\textbf{Architectural style.} Script-based repositories produce the highest mean improvement (+5.62) due to their few files and clear decomposition opportunities. Modular monoliths account for all five worsening cases, confirming that already-structured codebases are sensitive to structural interventions. Layered repositories show the highest improvement rate in count (6 of 9), suggesting that explicit layer boundaries provide a structured context for Localized Modification.

\textbf{Repository size.} Size is the dominant moderating variable. On tiny repositories, a single file split can raise the repository-level MI by over 20 points because the average is computed over only a few files. On large repositories (80+ files), individual changes are diluted by averaging and the median $|\Delta MI|$ is effectively zero. The correlation between Python file count and $|\Delta MI|$ is $-0.13$, confirming a weak but consistent inverse relationship.

\subsubsection*{Key Improvement Mechanism}

Across 12 of the 18 improved repositories, the core mechanism is the extraction of a new module achieving MI\,=\,100. These modules contain cohesive, simple code: configuration objects, pure utility functions, or API clients with limited branching. A secondary pattern involves files with MI\,=\,0 (Radon rank C): four repositories improved by at least 7 MI points after the LLM extracted part of a critically complex file, because the Halstead Volume component of MI decreases non-linearly with file size. Representative examples are:

\begin{itemize}[nosep]
    \item \textit{kodekloudhub/webapp-color} ($\Delta MI = +21.98$): a single \texttt{app.py} was split into \texttt{app.py}, \texttt{color.py} (MI\,=\,100), and \texttt{config.py} (MI\,=\,100).
    \item \textit{AutoLab-SAI-SJTU/Paper2Rebuttal} ($\Delta MI = +18.14$): \texttt{rebuttal\_service.py} (MI\,=\,0) was decomposed into three focused modules; the original file recovered to MI\,=\,55.6.
    \item \textit{HKUDS/FastCode} ($\Delta MI = +2.11$, 81 files): a rare large-repository success; \texttt{retriever.py} improved from MI\,=\,11.3 to 78.4 and \texttt{main.py} from 4.9 to 77.0 through targeted extraction.
\end{itemize}

\subsubsection*{Failure and Degradation Analysis}

Stable outcomes (17 cases) arise from three causes: generated files containing syntax errors that Radon excludes from the average (3 cases), planner failures in which the agent cannot produce a valid modification plan after repeated retries (7 cases), and no-op changes such as docstring additions that do not affect MI. Worsening outcomes (7 cases) occur when newly created files have MI below the existing repository average; all seven are bounded above $-1.24$ MI points and all occur in small or tiny repositories.

Technical pipeline failures include Ollama API rate limiting (HTTP 429), repeated planner JSON parse errors, stuck modification loops, and invalid step proposals. These failure modes indicate that multi-step agentic LLM pipelines require more robust planning validation and loop-termination strategies.


\section{Integrated Discussion}
\label{sec:eval_discussion}

The three studies together provide a multi-level view of what works and what is currently limited in LLM-based architecture-aware maintainability improvement.

\subsection{Architecture Detection Is Feasible but Boundary-Limited}

Study 2 shows that LLM-based architecture detection can produce above-baseline accuracy with a lightweight scriptable pipeline. No GPU or fine-tuning is required, and the full classification pipeline uses only AST-based extraction and standard cloud API access.

However, 70.2\% accuracy means approximately 30\% of repositories are misclassified. Most of these errors are modular monolith classified as layered. In the context of the pipeline, this error is particularly costly: a system that is domain-organized will receive layer-specific tactic recommendations designed for a technical layer hierarchy, which may be inappropriate or harmful.

This connects directly to the degradation cases observed in Study 1, where all seven degradations occurred in modular monolith repositories.

\subsection{Import Graphs Are the Most Valuable Architectural Signal}

The evidence type evaluation in Study 2 provides a clear design recommendation for any LLM-based architecture analysis system: include import dependency graphs, omit code signatures.

Import graphs encode structural dependencies between modules. They make visible the patterns that define architectural styles: whether dependencies flow in one direction (layered), whether modules are roughly symmetric (modular monolith), or whether there are few dependencies overall (script-based). This information is both inexpensive to extract and architecturally meaningful.

Code signatures, in contrast, added context volume without adding structural signal, and appear to have introduced misleading naming cues.

\subsection{LLM Transformations Remain Code-Level Rather Than Architecture-Level}

Studies 1 and 3 both show that the pipeline selects tactics at the architectural level but implements them through code-level modifications. Package count remained unchanged in all 123 repositories with paired metric data. Average fan-out changed in only 28 repositories, with a small mean change of approximately -0.016.

The Maintainability Index, which reflects file-level complexity through Halstead Volume and Cyclomatic Complexity, can capture these changes. However, architectural indicators such as module boundary structure, top-level package organization, and inter-module coupling patterns were largely unaffected.

This is consistent with the broader limitation observed in LLM-based refactoring research: even with repository-level context, LLMs tend to produce localized edits rather than restructuring. The implementation stage of the pipeline is therefore the primary bottleneck, independent of architecture detection quality.

\subsection{Script-Based Repositories Are the Strongest Use Case}

Across both Study 1 and Study 3, the strongest improvements occurred in script-based repositories where the Decomposability tactic was applied. These repositories typically consist of large monolithic files with high Halstead Volume and weak internal structure. Splitting such files into smaller cohesive modules directly reduces the inputs to the MI formula.

This represents a practical sweet spot: low-structure repositories offer clear decomposition targets, and even simple modularization produces measurable metric gains. The pipeline is most useful as a targeted tool for this type of repository rather than as a universal architecture-level refactoring system.

\subsection{The Confidence Problem Propagates Across Studies}

Study 2 demonstrated that self-reported LLM confidence does not distinguish correct from incorrect predictions. The maximum observed difference between correct and incorrect confidence was 0.026. Yet Study 1 relied on self-reported confidence as the sole indicator of architecture detection quality.

This means the interpretation of Study 1 architecture-level results must be treated cautiously. The tactic selection in Study 1 may have been partially based on incorrect architectural classifications, and confidence values cannot be used retroactively to identify which classifications were reliable.

Future iterations of the pipeline should use confidence scores only with significant caution, or replace them with structural consistency checks or ensemble agreement across models.

\subsection{Failure Rate Requires Verification Mechanisms}

The 25.3\% failure rate in Study 1 indicates that the current implementation mechanism is not reliable enough for autonomous deployment. Failures include syntax errors, broken imports, and structural changes that prevent re-analysis.

A production-ready implementation would require at minimum:
\begin{itemize}
    \item automated test suite execution before and after modification;
    \item build and import validation;
    \item rollback capability;
    \item human review for high-risk transformations.
\end{itemize}

The current framework provides syntactic validation but not behavioral preservation. This is a fundamental gap that limits the applicability of the pipeline to experimental contexts.

\section{Threats to Validity}
\label{sec:eval_threats}

\subsection{Internal Validity}

In Study 1, each repository serves as its own control, reducing cross-project variance. However, MI changes may be influenced by factors beyond the intended tactic, such as code formatting side effects or coincidental complexity changes during modification. A single LLM model was used without comparison to alternatives or manual implementation of tactics.

In Study 2, the ground truth labels were assigned by a single annotator. The modular monolith/layered boundary is inherently ambiguous, and a second annotator may draw the boundary differently for some repositories. Each prompt was executed once per model; the low temperature setting (0.2) reduces but does not eliminate stochastic variance.

\subsection{External Validity}

All three studies are limited to Python backend repositories from GitHub. The taxonomy covers three architectural styles; microservice, event-driven, and hexagonal architectures are not represented. Results may not generalize to other languages, enterprise systems, or more architecturally complex repositories.

\subsection{Construct Validity}

The Maintainability Index is a composite metric that does not capture all dimensions of architectural maintainability. It is sensitive to file-level complexity but less sensitive to modularity, architectural conformance, and inter-module dependency quality. Many meaningful modifications may not register as MI changes.

Architecture style is a high-level abstraction that some repositories may exhibit in hybrid form. Forcing repositories into a single class label may not reflect their actual structure.

\subsection{Conclusion Validity}

The datasets used in all three studies are relatively small. Study 2 uses 57 repositories; Study 1 uses 162 repositories with 121 completed. Statistical power is limited, and all results should be interpreted as early empirical evidence rather than definitive claims. All rankings and effect sizes are preliminary baselines.
